\subsection{Cube-chordal ample classes}
\label{sec:cube-chordal}

Chordality gives another sufficient condition for corner peeling, expressed
in terms of the actual maximal cubes of a class. We collect its definition,
structural descriptions and closure properties here. We introduce the name
\emph{cube-chordal ample} for this class; \emph{daisy cube} is established
terminology \cite{KlavzarMollardDaisy}.
The relative peeling theorem is recalled from the dismantlable-ample
companion, and the tree description from the maximum-ancestry companion.
The product criterion and explicit examples make their consequences precise.
The complementary completion theorem below also puts this class in
$\mathcal B$, where both original ambient shores are corner-peelable.

\begin{definition}[Cube-chordal ample classes]
\label{def:cube-chordal}
For nonempty $H\subseteq\bits^E$, a positioned cube is
\[
 C(p,S)=\{x\in\bits^E:x|_{E\setminus S}=p\},
 \qquad S\subseteq E,\quad p\in\bits^{E\setminus S}.
\]
Its free coordinate set is $S$. Let $\mathcal M(H)$ be the family of
\emph{all} inclusion-maximal positioned cubes contained in $H$.
The maximal-cube intersection graph $\Gamma(H)$ has vertex set
$\mathcal M(H)$; distinct cubes are adjacent exactly when their vertex
sets intersect, including an intersection consisting of one word.
We call $H$ \emph{cube-chordal ample} if $H$ is ample and $\Gamma(H)$
is chordal, meaning that it has no induced cycle of length at least four.
Write $\mathsf{CAmp}$ for this family, and use \emph{cube-chordal} as
shorthand when ampleness is understood.
\end{definition}

The graph uses complete positioned cubes, rather than just their free
coordinate sets or a selected cube cover. Equivalently its vertices are
the extentures of $\bits^E\setminus H$: two such minimal forbidden
patterns are adjacent when their fixed literals are compatible.
Chordality here concerns $\Gamma(H)$, not the one-inclusion graph of $H$;
even a full square has a nonchordal one-inclusion graph.

A \emph{daisy cube rooted at $z$} means that
$\{\operatorname{Diff}(z,x):x\in H\}$ is a downset in $2^E$,
where $\operatorname{Diff}(z,x)=\{i\in E:z_i\ne x_i\}$.
Equivalently, $H$ contains the full cube interval from $z$ to each of
its words: $I(z,x)=\{y:\operatorname{Diff}(z,y)\subseteq
\operatorname{Diff}(z,x)\}\subseteq H$ for every $x\in H$.
Thus the root specifies the Boolean orientation of the established daisy
cube notion recalled above. These classes are ample: if $S$ is shattered,
a word differing from $z$ on all of $S$ supplies the entire rooted
$S$-cube by downward closure.

\begin{proposition}[The complete-graph case]
\label{prop:cube-chordal-daisy}
For a nonempty class $H$, $\Gamma(H)$ is complete if and only if $H$ is
a daisy cube. Consequently
\[
 \mathsf{Daisy}\subsetneq\mathsf{CAmp}\subsetneq\Damp\subseteq\Rep.
\]
There is no bound on VC dimension in $\mathsf{CAmp}$.
\end{proposition}
\begin{proof}
Positioned cubes have the Helly property: pairwise intersection says that
no two fixed literals conflict, so all the fixed literals are jointly
consistent. If $\Gamma(H)$ is complete, all maximal cubes therefore
contain a common word $z$. Each is downward closed towards $z$, and so
is their union $H$. Conversely the maximal cubes of a daisy cube all
contain its root: the top word of any contained cube supplies its whole
rooted lower cube. Thus their intersection graph is complete.

The four-word path $000,100,110,111$ has three maximal edges and
$\Gamma(H)=P_3$. It is ample, since its three directions are distinct,
and is cube-chordal but not a daisy cube. The inclusion into $\Damp$
is proved below, and its strictness by
\cref{ex:cube-chordal-corner}. The last inclusion is the existing
corner-peeling representation theorem. Full cubes give unbounded VC dimension.
\end{proof}

\subsubsection{Daisy pieces, cube attachments and prescribed endpoints}

The following gives an exact construction in the original Boolean
coordinates. It uses the complete extenture family, so no choice of a
redundant forbidden-pattern presentation is implicit.

\begin{proposition}[An incremental construction]
\label{prop:cube-chordal-generator}
Start from any one positioned cube. Given a previously constructed class
$B$, attach $C=C(\delta,S)$, with $Y=E\setminus S$, subject to:
\begin{enumerate}[label=(\roman*)]
\item $C\not\subseteq B$, and $C$ contains no member of $\mathcal M(B)$;
\item the clamp $c(b)=(b|_S,\delta|_Y)$ belongs to $B$ for every $b\in B$;
\item the members of $\mathcal M(B)$ meeting $C$ intersect pairwise.
\end{enumerate}
The outputs are exactly $\mathsf{CAmp}$. Condition~(ii) is equivalent to
\[
 \pi_SB=\pi_S(B\cap C),
\]
and to the following extenture condition: every extenture of $B$ whose
support meets $Y$ has a $Y$-literal opposite to $\delta$.
The interface $F=B\cap C$ is a nonempty proper daisy cube inside $C$,
rooted at a common word of $C$ and all its neighbours.
It is gated in $B$, with gate $c(b)$: for every $b\in B$ and $f\in F$,
\[
 d(b,f)=d(b,c(b))+d(c(b),f).
\]
Here $d$ is Hamming distance, equal to graph distance on the ample pieces.
\end{proposition}
\begin{proof}
First suppose $B$ is ample and (ii) holds. Every cube in $B\cup C$
lies in one piece. Otherwise it contains a new word $v\in C\setminus B$
and frees some $Y$-coordinate; flipping that coordinate at $v$ gives a
word $b$ outside $C$, hence in $B$, with $c(b)=v$, a contradiction.
Also every support shattered by $B\cup C$ that meets $Y$ is already
shattered by $B$. To realize a prescribed trace whose $Y$-part agrees
with $\delta$, first flip one of its prescribed $Y$-bits and realize that
trace in the union. Its witness lies in $B$; clamping gives the desired
trace in $B$. Traces not agreeing with $\delta$ already have witnesses
in $B$. Thus
\[
 \Sh(B\cup C)=\Sh(B)\cup2^S.
\]
The same union describes strongly shattered supports by the cube assertion,
so $B\cup C$ is ample. By (i) its maximal cubes are exactly
$\mathcal M(B)\cup\{C\}$. Condition~(iii) makes the new graph vertex
simplicial, preserving chordality. This proves soundness by induction.

For the converse, use the \href{../maxmin-forbidden/sections/chordal_guard_coefficients.tex}
{simplicial-cube remainder lemma} in the maximum-ancestry companion
(label \nolinkurl{lem:maxmin-simplicial-cube-remainder}).
It states that removing the exclusive words of a simplicial maximal cube
from an ample class leaves an ample $B$ with exactly the other maximal
cubes. Here is the removal argument: cube Helly supplies a common word
$z$ of that cube and its neighbours. Their intersections form a proper
daisy cube $F$ rooted at $z$. Delete the remaining words in decreasing
support inclusion relative to $z$. Each is a corner, because it is maximal
in the remaining downset inside $C$ and belongs to no other old cube.
Every cube of $F$ lies in one neighbouring cube: its top word and $z$
belong to that neighbour's face in $C$. This also proves the exact remaining
maximal-cube list.

For this remainder every cube of $B\cup C$ lies in a piece. Ampleness and
the Sandwich equality $|A|=|\Sh(A)|$ for ample $A$ therefore give
\[
 |\Sh(B)\cap2^S|=|B|+|C|-|B\cup C|=|F|.
\]
Projection preserves ampleness, so $|\pi_SB|=|\Sh(B)\cap2^S|$.
Since $\pi_SF\subseteq\pi_SB$ has $|F|$ words, the projections agree.
This proves (ii); maximality and simpliciality give (i),(iii).
A chordal graph has a simplicial vertex, and deleting it leaves a chordal
graph. Repeat until one maximal cube remains, then reverse the steps.

For the extenture equivalence, an extenture meeting $Y$ with every
$Y$-literal equal to $\delta$ would, by minimality, have its $S$-part
realized in $B$. Clamping that witness would realize the forbidden pattern.
Conversely a clamp avoids every $S$-only extenture because its original
word did, and avoids every other extenture by an opposite $Y$-literal.
Helly and (iii) give the asserted common root and the daisy interface.
It is nonempty by (ii), and proper by (i). Finally $b$ and $c(b)$ differ
only on $Y$, whereas $c(b)$ and $f$ differ only on $S$; this proves the
distance identity and gatedness. Uniqueness follows by putting $f=c(b)$
in the corresponding identity for any proposed gate.
\end{proof}

The interface may be a union of faces rather than a single cube.
For example, attach the three maximal squares of
$\bits^3\setminus\{111\}$ in succession. The last interface is a
three-word path, a daisy cube with two maximal edges. It is gated in the
old host; it need not be gated in the incoming square.

There is also an exact simultaneous tree description. The maximal cliques
of $\Gamma(H)$ have a clique tree: every cube label occurs in a connected
subtree, and intersecting labels occur together in a bag. Each bag union
is a daisy cube rooted at its common word. Adjacent bag unions intersect
in the union of their shared cube labels, a daisy cube rooted at both bag
roots. These assertions, including a converse, are proved in
\href{../maxmin-forbidden/sections/chordal_downset_tree.tex}
{\nolinkurl{thm:maxmin-chordal-downset-tree}}.
The converse requires one additional compatibility condition. Write $d(C)$
for the free coordinate set of a positioned cube $C$. For every
coordinate face $P$ and $S\subseteq d(P)$, the union of occurrence subtrees
of labels $C$ with $C\cap P\ne\varnothing$ and $S\subseteq d(C)$ must
be connected or empty. These are unions of tree vertices \emph{and edges};
adjacent vertices alone do not supply an edge. An arbitrary tree of daisy
cubes need not satisfy this condition and is not an equivalent definition.

\begin{proposition}[Peeling to any contained cube]
\label{prop:cube-chordal-peeling}
If $H\in\mathsf{CAmp}$ and $Q\subseteq H$ is a nonempty positioned
cube, there is a corner order from $H$ to $Q$ in which every remaining
class is cube-chordal ample. Every supplied ordinary representation on $Q$
extends to $H$, with boundary edges directed into $Q$ and no directed
cycle outside $Q$. In particular every vertex can be the final word of
a corner order and the sink of an acyclic ordinary representation of $H$.
\end{proposition}
\begin{proof}
This is the cube-retention theorem
\cite[\nolinkurl{thm:damp-chordal-cube-retention}]{AmpleCampaign}.
We recall the argument to specify the strength of the statement.
A chordal graph has a simplicial vertex outside any proper clique:
a noncomplete chordal graph has two nonadjacent simplicial vertices,
and every vertex of a complete graph is simplicial. Protect the clique
of maximal cubes containing $Q$. If it is proper, choose a simplicial
cube $C$ outside it. Let $z$ be a common word of $C$ and its neighbours,
and $y$ its antipode in $C$. Every neighbouring face contains $z$ and
excludes $y$, so $y$ has unique maximal cube $C$ and is a corner outside $Q$.
If all maximal cubes contain $Q$ and $H\ne Q$, choose any such $C$,
take $z\in Q$, and again its antipode $y\notin Q$ is a corner.

Deleting $y$ replaces $C$ by its facets opposite $y$, discarding facets
contained in unchanged cubes. The retained facets form a clique with
exactly the old clique neighbourhood of $C$, since they all contain $z$.
Each is simplicial, so chordality is preserved. Corner deletion preserves
ampleness. Induction reaches $Q$.
Reverse the order, adding each removed corner as a source. This extends
an ordinary representation by the corner-extension theorem, preserves the
map on $Q$, and makes its boundary inward. A directed cycle cannot use
an added vertex: the last added vertex on that cycle would have only
outgoing edges to earlier vertices. For a prescribed single-vertex target,
its trivial map therefore extends acyclically with that vertex as sink.
\end{proof}

\begin{samepage}
\subsubsection{Corner peelings on both ambient shores}

\begin{proposition}[Complementary completion of cube-chordal classes]
\label{prop:atlas-cube-chordal-complement}
If $H\subseteq\bits^E$ is nonempty and cube-chordal ample, then both $H$
and its full original ambient complement have corner peelings. Consequently
\[
 \mathsf{CAmp}\subseteq\mathcal B\subsetneq\mathsf{Ext}.
\]
There is no bound on VC dimension. Original constant coordinates are retained;
an empty complement has the empty map and peeling order.
\end{proposition}
\end{samepage}
\begin{proof}
Apply the tree completion theorem
\cite[\texttt{prop:leaf-root-alignment}, \texttt{thm:camp-ext}]{CubeComplementCompletion}.
Its input is the positioned tree of daisy pieces recalled above: every
connected subunion is ample, every cube is carried by a piece, word
occurrences are connected, and each adjacent overlap is rooted at both
piece roots. The theorem supplies a finite increasing sequence
$H=H_0\subset\cdots\subset H_m=L$, where $L$ is one daisy cube.
Each step attaches a cube $Q(z,S)$ whose clamp preserves $H_j$ and whose
intersection with $H_j$ is a nonempty daisy rooted at $z$.
The construction preserves the tree properties; it does not assume that
intermediate maximal-cube graphs remain chordal.

For completeness, such a step gives an explicit complementary peeling.
The added words $A=Q(z,S)\setminus H_j$ form an up-set towards the corner
opposite $z$. Remove them from $H_j^c$ in increasing distance from $z$.
At the removal of $a$, put $D=\operatorname{Diff}(a,z)\subseteq S$.
Rootward neighbors are old words or were removed earlier. The entire
cube with directions $E\setminus D$ remains: its traces in $Q(z,S)$
are later words of $A$, and all parallel outside words are absent from
$H_j$ by clamp invariance. Thus $a$ is an actual corner with that support.
The remaining class is $H_{j+1}^c$. After all steps, peel the up-set $L^c$
in increasing distance from its daisy root. The primal peeling is
\cref{prop:cube-chordal-peeling}. This proves membership in $\mathcal B$;
$\mathcal B\subsetneq\mathsf{Ext}$ was established in
\cref{prop:maxmin-neighbors}.
\end{proof}

Neither complementation nor projection must preserve cube chordality.
The conclusion above concerns the two shores' corner peelings, not their
maximal-cube graphs. It also does not supply maximum parents:
$\mathsf{CAmp}\subseteq\MaxMin$ remains open.
Conversely, if an ample $H$ is not corner-peelable, then neither $H$ nor
$H^c$ is cube-chordal. Thus both maximal-cube intersection graphs contain
an induced cycle of length at least four; this is a necessary condition,
not a characterization of nondismantlability.

\subsubsection{Closure properties and the exact product criterion}

In the table, empty results may be adjoined formally. A full reduction on
$I$ keeps exactly those words whose entire $I$-fiber is present:
\[
 H^I=\{x\in\bits^{E\setminus I}:\{x\}\times\bits^I\subseteq H\}.
\]
A rooted vertex wedge identifies specified vertices of two classes on
disjoint coordinate blocks, keeping all coordinates of the other block
fixed at that vertex in each piece.

\begin{center}
\begin{tabular}{@{}p{.72\linewidth}l@{}}
\toprule
Operation & Preserves $\mathsf{CAmp}$?\\
\midrule
Signed coordinate isometries; adding or suppressing constants & Yes\\
Coordinate faces; geodesically convex subclasses & Yes\\
Full coordinate reductions $H^I$ & Yes\\
Rooted vertex wedges on disjoint blocks & Yes\\
Product with any daisy cube & Yes\\
Arbitrary Cartesian products & No\\
Coordinate projections; arbitrary pc-minors & No\\
Complementation in the specified ambient cube & No\\
Arbitrary corner deletion or corner addition & No\\
Gluing along a common cube, even a convex coordinate face & No\\
Arbitrary ample subclasses & No\\
\bottomrule
\end{tabular}
\end{center}

\begin{proposition}[Faces, reductions and vertex wedges]
\label{prop:cube-chordal-closures}
The affirmative closure assertions in the table hold.
\end{proposition}
\begin{proof}
Signed isometries and constants preserve the maximal-cube graph.
For a coordinate face $P$, its maximal cubes in $H\cap P$ are the maximal
nonempty traces $C\cap P$ of old maximal cubes. Choose one old cube for
each such trace. Two retained traces intersect if and only if their old
representatives intersect: the reverse implication uses cube Helly on
the two representatives and $P$. Thus $\Gamma(H\cap P)$ is isomorphic
to an induced subgraph of $\Gamma(H)$ and is chordal. Faces preserve
ampleness. Convex subclasses of a partial cube are intersections of its
coordinate halfspaces, as recalled in the operations section, so are faces.

For reduction, the maximal cubes of $H^I$ correspond exactly to the old
maximal cubes freeing every coordinate of $I$, with $I$ removed.
Indeed lift a reduction cube to its product with $\bits^I$ and extend it
to a maximal cube of $H$; maximality on either side gives the converse.
Intersections are preserved, so the graph is again induced. Full reductions
preserve ampleness by iterating the usual ample two-sided reduction.

Singleton wedge factors are neutral. Otherwise all maximal cubes come
from one factor. Cubes from different factors meet exactly when both
contain the identified root. The root-containing cubes of each factor
form a clique. Glue the two chordal graphs to the complete graph on
their combined root-containing cubes, along these two cliques.
Clique gluing preserves chordality: an induced cycle visiting both sides
of a clique separator would have a chord. Finally a shattered support
cannot use a nonconstant coordinate from each private block, since the
pattern opposite both roots is absent. Supports wholly in one block are
shattered exactly when that factor shatters them, and are strongly
shattered there. Thus the wedge is ample. Product closure with daisy cubes
is the next proposition.
\end{proof}

\begin{proposition}[Exact product criterion]
\label{prop:cube-chordal-product}
For any nonempty ample $H,K$ on disjoint coordinate sets,
\[
 H\times K\in\mathsf{CAmp}
 \quad\Longleftrightarrow\quad
 H,K\in\mathsf{CAmp}\ \text{and at least one factor is a daisy cube}.
\]
\end{proposition}
\begin{proof}
Products of ample classes are ample: a support is shattered, or strongly
shattered, exactly when its two factor supports are. The maximal cubes
of a product are precisely the products of maximal cubes. Consequently
$\Gamma(H\times K)=\Gamma(H)\boxtimes\Gamma(K)$, the strong graph
product: two distinct pairs are adjacent when, in each coordinate, their
entries are equal or adjacent. Fixing either graph entry gives the other
graph as an induced subgraph, so both graphs must be chordal.

These graphs are connected. Indeed an ample class is connected, and
its vertex paths give overlap chains of maximal cubes. Any connected
noncomplete graph contains an induced path on three vertices: use the
first three vertices of a shortest path between nonadjacent vertices.
If both factor graphs are noncomplete, choose such paths indexed $0,1,2$.
The four pairs
\[
 (0,1),\ (1,0),\ (2,1),\ (1,2)
\]
form an induced four-cycle in their strong product. Thus at least one
graph must be complete. Conversely, a strong product with a complete
graph replaces each vertex of the other graph by a clique, giving all
vertices in that clique the same outside neighbours. Replacing each
vertex in a perfect elimination order by its whole clique gives a
perfect elimination order of the product. Chordality is preserved.
Apply \cref{prop:cube-chordal-daisy} to identify the complete factor.
\end{proof}

For instance the four-word path in
\cref{prop:cube-chordal-daisy} is cube-chordal, whereas its product
with itself is not. This product is median, hence in $\Rper$.
Thus even median or recursively peripheral membership does not imply
cube-chordality.

\subsubsection{Explicit failures of closure}

For compact examples, encode $x\in\bits^{\{0,\ldots,n-1\}}$ by
$\sum_i2^ix_i$. Write $Q(a,s)$ for the cube whose free coordinates
are the one-bits of $s$ and whose fixed one-bits are those of $a$;
the two masks have disjoint one-bits. Thus its words are $a\mathbin{\mathrm{or}}t$
for the submasks $t$ of $s$. All lists below are complete maximal-cube lists.

\begin{example}[Corner changes and common-cube gluing]
\label{ex:cube-chordal-corner}
On four coordinates let
\[
 A=\{0,1,2,3,6,7,8,9,10\},\qquad D=A\cup\{5\},\qquad
 U=D\cup\{4\}=\{0,1,\ldots,10\}.
\]
Their maximal cubes, listed as $(a,s)$, are
\begin{align*}
 \mathcal M(A)&:\ (0,3),(2,5),(0,9),(0,10),\\
 \mathcal M(D)&:\ (0,3),(2,5),(1,6),(0,9),(0,10),\\
 \mathcal M(U)&:\ (0,7),(0,9),(0,10).
\end{align*}
These lists follow by extending each edge and face in the displayed
word sets. In $D$, the four cubes
$(2,5),(1,6),(0,9),(0,10)$ form an induced four-cycle in that order;
$(0,3)$ meets all four. Thus $\Gamma(D)$ is the wheel on five vertices,
whereas $\Gamma(A)$ is a diamond and $\Gamma(U)$ is complete.

The class $U$ is a daisy cube rooted at zero. The word $4$ is its corner
with incident free mask $7$, so $D$ is ample and dismantlable.
The word $5$ is a corner of $D$ with incident free mask $6$, leaving $A$.
Hence $U\mapsto D$ disproves closure under corner deletion, and
$A\mapsto D$ disproves closure under corner addition. The inclusion
$D\subseteq\bits^4$ also disproves closure under ample subclasses.
Since $D\in\Damp\setminus\mathsf{CAmp}$, the inclusion into $\Damp$
is strict.

Moreover $D$ is the union of the two daisy cubes
\[
 Z=\{0,1,2,3,5,6,7\},\qquad W=\{0,1,2,3,8,9,10\},
\]
rooted at $3$ and $0$, respectively. Their intersection is the full
square $\{0,1,2,3\}$, a coordinate face and therefore convex in each.
This disproves common-cube gluing closure even for daisy input pieces.
\end{example}

\begin{example}[Projection and complementation]
\label{ex:cube-chordal-projection}
On five coordinates put
\[
 V=\{0,1,2,3,6,7,8,9,16,17,18,19,21,23,24,26\}.
\]
Its maximal cubes are
\[
 (0,19),(2,5),(17,6),(0,9),(16,10),(3,20),(0,24).
\]
Their graph is the cone over two disjoint three-vertex paths and is
chordal. Ampleness can be verified directly. At $x_2=x_3=0$ there is a
full cube on coordinates $\{0,1,4\}$. At $x_2=1$ there is a four-word
path on these coordinates using each direction once, and the same holds
at $x_3=1$. No word has $x_2=x_3=1$. Consequently the shattered supports
are the eight subsets of $\{0,1,4\}$, together with
$\{\{j\}\}\cup\{\{j,i\}:i\in\{0,1,4\}\}$ for each $j=2,3$.
Each is strongly shattered by the central cube, a vertical edge or a
path-edge square. Hence $V$ is ample. Projecting away coordinate $4$
gives exactly $D$ of \cref{ex:cube-chordal-corner}, proving failure
of projection closure and therefore pc-minor closure.

For complementation take the four-coordinate path
$P=\{0,1,3,7,8\}$. Its four edge directions are distinct, so it is
maximum of VC dimension one, with a path as maximal-cube graph.
Its complement is ample by ample complement duality and has maximal cubes
\[
 (12,3),(10,5),(9,6),(4,9),(4,10),(2,12).
\]
The cubes $(10,5),(9,6),(4,9),(4,10)$ form an induced four-cycle in
that order. This proves failure of complement closure.
\end{example}

\subsubsection{Comparison with ACC and maximum ancestry}

Recall that an ACC class has a join tree on its individual maximal
cubes, as defined before \cref{fig:classes}; see also the
\href{../../../YangDim/yang_complex_dictionary.tex}{Yang-complex dictionary}.
A clique tree of $\Gamma(H)$ instead has \emph{bags} of cubes as nodes.
These two requirements are incomparable.

\begin{example}[Neither ACC nor medianity characterizes cube-chordality]
\label{ex:cube-chordal-acc}
The daisy cube $\bits^3\setminus\{111\}$ is cube-chordal, but its three
maximal squares admit no join tree. Each pair has a word belonging to
just those two squares, forcing each of the three possible tree edges.
Even medianity does not repair the implication: the $3$-by-$3$ grid is
a median daisy cube with four maximal squares sharing its middle vertex,
so $\Gamma=K_4$. Each of the four neighbours of that vertex belongs
to just one cyclically consecutive pair of squares. Those four pairs
would all have to be edges of a join tree, which is impossible.

Conversely take a central square and attach a square along each of its
four boundary edges, using a separate private coordinate for each leaf.
The central square and its four leaves have a join tree given by the star.
The class is isometric: within a leaf use square paths; between different
leaves, turn off their distinct private coordinates and take a shortest
path in the central square. It is also ample: its shattered supports are
the empty set, six singletons and the five pairs supporting its squares,
all strongly shattered. Thus it is ACC. Its four leaf squares meet in
an induced four-cycle, so its maximal-cube graph is nonchordal.
\end{example}

\begin{proposition}[Minor closure of ACC]
\label{prop:atlas-acc-closure}
The ACC classes are closed under nonempty projections and conditionings.
Their inclusion into Median is strict.
\end{proposition}
\begin{proof}
Use the characterizations in
\cite[Theorem~\textnormal{ACC characterizations}]{Survey}.
For a median class the split-incompatibility graph has one vertex for
each nonconstant coordinate, joining a pair exactly when it is shattered.
Projection retains its induced graph on the remaining active coordinates.
Chordality and medianity are preserved, proving projection closure.

For conditioning, let $G^\Delta$ join two concepts whenever a contained
cube contains both. If $F$ is a coordinate face, intersecting an old cube
with $F$ shows that $(G\cap F)^\Delta$ is the induced subgraph
$G^\Delta[V(G)\cap F]$. An ACC class has chordal $G^\Delta$;
its induced subgraph is chordal as well. The nonempty coordinate face
is median, so the same characterization proves it ACC.
The median $3$-by-$3$ grid in \cref{ex:cube-chordal-acc} is not ACC,
proving strictness.
\end{proof}

We have proved
$\mathsf{Daisy}\subsetneq\mathsf{CAmp}\subsetneq\Damp$.
The product of the four-word path with itself proves
$\Rper\not\subseteq\mathsf{CAmp}$.
For the other direction the peripheral companion's
\href{../cover-recursive/research/rper_chordal_sun_obstructions.tex}
{chordal sun family} has chordal maximal-cube graph and is forbidden for
$\Rper$. Hence $\mathsf{CAmp}\not\subseteq\Rper$ as well.
Those suns do have maximum parents, as shown in
\href{../maxmin-forbidden/research/chordal_sun_parents.tex}
{\nolinkurl{thm:maxmin-chordal-sun-parents}}.
The unrestricted inclusion $\mathsf{CAmp}\subseteq\MaxMin$ remains open;
neither a clique tree nor the local generator supplies compatible maximum
parents for arbitrary daisy interfaces. The class is not pc-minor closed,
so a forbidden-pc-minor basis cannot characterize membership in
$\mathsf{CAmp}$ itself.
